\section{Factor-quotient jets and Pell projective curvature}
\label{sec:pell-factor-quotient-projective}

Write
\[
 (1+\sqrt2)^n=A_n+B_n\sqrt2,\qquad
 U=A_\ell B_\ell,\qquad v=2A_{2\ell},\qquad x=2\ell,
\]
where $\ell$ is an odd prime.  Put
\[
 a=\frac{A_\ell-1}{x},\qquad
 b=\frac{s_\ell B_\ell-1}{x},\qquad
 s_\ell=\left(\frac2\ell\right).
\]
For a prime factor $q\mid A_\ell$, write $q=1+xk_q$ and repeat $k_q$
according to its valuation.  For $r\mid B_\ell$, write
$r=s_r+xh_r=s_r(1+xs_rh_r)$ and repeat $s_rh_r$.  If the first three
elementary symmetric functions of the two resulting lists are
$(K_A,C_A,H_A)$ and $(K_B,C_B,H_B)$, respectively, exact finite-product
expansion gives
\[
 a\equiv K_A+xC_A+x^2H_A\pmod{x^3},\qquad
 b\equiv K_B+xC_B+x^2H_B\pmod{x^3}.                    \tag{\ref{sec:pell-factor-quotient-projective}.1}
\]
Three copies of one quotient $t$ have the exact fingerprint
\[
\begin{aligned}
 E_1(t,t,t,R)&=3t+K_0,\\
 E_2(t,t,t,R)&=3t^2+3tK_0+C_0,\\
 E_3(t,t,t,R)&=t^3+3t^2K_0+3tC_0+H_0,
\end{aligned}                                           \tag{\ref{sec:pell-factor-quotient-projective}.2}
\]
where $(K_0,C_0,H_0)$ is the residual triple.  A depth-three carrier is
therefore visible in the third-order ledger even when its actual depth is
larger.

The factor quotients also determine a jet of the Lucas companion.

\begin{theorem}[Two-channel companion jet]
\label{thm:pell-two-channel-companion-jet}
Define
\[
\begin{aligned}
 V_A(K,C,H)&=6+8xK+x^2(8C+4K^2)+x^3(8H+8KC),\\
 V_B(K,C,H)&=6+16xK+x^2(16C+8K^2)+x^3(16H+16KC).
\end{aligned}
\]
Then
\[
 v\equiv V_A(K_A,C_A,H_A)\pmod{x^4},\qquad
 v\equiv V_B(K_B,C_B,H_B)\pmod{x^4}.                  \tag{\ref{thm:pell-two-channel-companion-jet}.1}
\]
\end{theorem}

\begin{proof}
The negative Pell identity gives
$v=4A_\ell^2+2=8B_\ell^2-2$.  Substitute
$A_\ell=1+xa$ in the first expression and
$s_\ell B_\ell=1+xb$ in the second.  Modulo $x^4$, the square of
$K+xC+x^2H$ contributes only $K^2+2xKC$ after multiplication by $x^2$.
The displayed polynomials follow by collecting coefficients.
\end{proof}

Let $\theta=(\ell-1)/2$ and let $\alpha_j,\beta_j$ be the correlated
Lucas coefficients from the preceding section, satisfying
\[
 (2j+1)\alpha_j=\ell\beta_j,\qquad
 \alpha_\theta=\beta_\theta=32^\theta.
\]
For the tail polynomials $E_j,F_j$, put $T_j=vF_j(U^2)$ and define
\[
 \Delta_{\rm top}=
 T_{\theta-1}\ell E_\theta(U^2)
 -T_\theta(\ell-2)E_{\theta-1}(U^2).
\]

\begin{theorem}[Exact endpoint curvature]
\label{thm:pell-exact-endpoint-curvature}
For every odd prime index,
\[
                 \Delta_{\rm top}=2v\,32^{\ell-1}U^2.                 \tag{\ref{thm:pell-exact-endpoint-curvature}.1}
\]
Moreover
\[
 \gcd(2v32^{\ell-1},U)=1,\qquad
 v_p(\Delta_{\rm top})=2v_p(U)\quad(p\mid U).           \tag{\ref{thm:pell-exact-endpoint-curvature}.2}
\]
Thus $U^2$ always divides the determinant and $U^3$ never does.
\end{theorem}

\begin{proof}
The two top tails are
\[
 E_{\theta-1}=\alpha_{\theta-1}+\alpha_\theta U^2,
 \quad F_{\theta-1}=\beta_{\theta-1}+\alpha_\theta U^2,
 \quad E_\theta=F_\theta=\alpha_\theta.
\]
After substituting $T_j=vF_j$, the constant part of the determinant is
\[
 v\alpha_\theta
 \{\ell\beta_{\theta-1}-(\ell-2)\alpha_{\theta-1}\}=0
\]
by coefficient correlation.  Its $U^2$ coefficient is
$2v\alpha_\theta^2=2v32^{\ell-1}$, proving the first identity.
Odd-index Pell coordinates make $U$ odd, while
\[
                         v^2-32U^2=4.
\]
Every common odd divisor of $v$ and $U$ would divide four.  The quotient in
the first identity is therefore a unit at every prime of $U$, proving the
valuation statement and sharpness.
\end{proof}

With $\kappa_\ell=\Delta_{\rm top}/U^2=2v32^{\ell-1}$,
Theorems~\ref{thm:pell-two-channel-companion-jet} and
\ref{thm:pell-exact-endpoint-curvature} give the direct coupling
\[
 \kappa_\ell\equiv
 2\,32^{\ell-1}V_A(K_A,C_A,H_A)
 \equiv2\,32^{\ell-1}V_B(K_B,C_B,H_B)\pmod{16\ell^4}.  \tag{\ref{thm:pell-exact-endpoint-curvature}.3}
\]
Thus the third-order factor-quotient ledger determines the
$2\ell$-adic jet of the exact unit measuring failure of integral projective
collinearity.  Since the rank theorem gives $\gcd(2\ell,U)=1$, the jet
modulus and $U^2$ are coprime; the Chinese remainder theorem combines them
without contradiction.  A closing argument must relate their values through
reciprocity, height, or distribution.

Two full-premise counterexamples delimit stronger local claims.  First, the
local-only assertion that the forced depth-three quotient data and complete
third-order ledger are inconsistent is false.  Take
\[
 \ell=3,\qquad q=7=1+6,\qquad r=797=-1+6\cdot133,
\]
and quotient lists $[1,1,1]$ and $[-133,-133,-133]$.  They satisfy the
required congruence and character table, and the full ledger is divisible by
$8\ell^3$.  The synthetic values are squarefull of depth three.  They do
not satisfy the global Pell equation:
\[
 q^6-2r^6+1=-512601560592751008\ne0.
\]
This retires only the local inconsistency claim.

Second, the genuine Pell index $\ell=7$ has
\[
 A_7=239,\quad B_7=169,\quad U=40391,\quad v=228486,
\]
and
\[
 \Delta_{\rm top}=800495088185894730989568,\qquad
 \Delta_{\rm top}\bmod U^3=24354030047568\ne0.
\]
It is a complete counterexample to the proposed $U^3$ strengthening and an
instance of the exact $U^2$ theorem.  Since $239\parallel A_7$, it is not a
counterexample to a hypothetical squarefull Pell packet.

Independent certificates exhibit a prime divisor of exponent one in
$A_\ell B_\ell$ for every one of the $57$ odd prime indices
$3\le\ell\le271$.  This proves finite nonsquarefullness in that range and
nothing beyond it.  The active route must still rule out or construct an
unbounded family satisfying simultaneously squarefullness, the opposite
depth-three nonresidue pair, both factor-quotient jets, the complete
third-order ledger, all correlated Lucas tails, exact endpoint curvature,
and vertexwise character incidence.  No full-premise counterexample to that
global statement is known.
